\section{The construction and the results recovered}
We study the construction of Leemans [L], which extends an earlier
construction of Van Nypelseer. The source is the complete fourteen-page
journal edition, \emph{A generalization of a construction due to Van
Nypelseer}, Beitr\"age zur Algebra und Geometrie 46 (2005), 561--574.
For a Steiner design the construction is precisely the geometry of
Part II of the common paper [G]. For an arbitrary partial linear space
its rank-three form is given below. We call a firm, residually connected,
flag-transitive geometry satisfying $(\mathrm{IP})_2$ an A-geometry.
Automorphisms preserve types unless stated otherwise.

The purpose is to recover the principal classification and its geometric
and group data, while correcting false intermediate assertions. The
common theorems give all chamber and residue reductions. The remaining
work consists of explicit rank-three geometries, field actions, and
finite Witt-design stabilizers. In particular, residual connectedness
and flag transitivity are separate conditions: the second iteration of
a deleted affine line is an A-geometry at $q=8$, rather than the source's
$q=4$; the subline construction is residually connected for every $q\geq3$
when the extension degree is two, although its flag-transitive case still
has $q=4$.

The full Hermitian-unital proof is included, using the transport theorem
of [G, Part III]. Its conjectural statement in [L, \S4.6] was already
resolved in the earlier accepted project work; this edition does not
claim a second result or a new priority. The rank-six Mathieu geometry,
the distinct rank-five residues and their actual stabilizer groups are
also included. Identifications with Pasini's named geometries
$\Gamma_{23}(4)$ and $\Gamma_8(4)$ remain excluded: the available evidence
does not supply object-level maps to those earlier definitions. No
isomorphism of geometries is inferred merely from matching group names.

\subsection{Two outside classification theorems}
Two substantial classifications are used at their point of application;
they are not consequences of the elementary arguments in [G].
Buekenhout, Dehaye and Leemans [BDL, Theorem 6.2, pp.42--43] classify
finite linear spaces with a faithful flag-transitive group whose point
and line stabilizers are doubly transitive on their respective incident
objects. Their possibilities are:
\begin{enumerate}
\item $AG(2,4)$ with $A\Gamma L(1,16)$;
\item $PG(n,q)$, $n\geq2$, with
$PSL(n+1,q)\leq G\leq P\Gamma L(n+1,q)$;
\item the specified $A_7$ action on $PG(3,2)$;
\item $AG(n,q)$, $n\geq2$, $q\geq3$, with translations and linear
stabilizer containing $SL(n,q)$;
\item the classical unital $U_H(q)$, with
$PSU(3,q)\leq G\leq P\Gamma U(3,q)$;
\item circles, with the specified triply transitive point actions.
\end{enumerate}
The last case is absent under the thickness hypothesis of [L, Theorem 2].
Small overlapping cases are not separate families; for example
$U_H(2)\cong AG(2,3)$. This input supplies the classification quoted in
that source theorem. The construction-specific conclusions for the
projective and affine cases will be proved directly below.

For the higher-rank classification we use Kantor [K, Theorem 3 and its
proof, p.73]. For a Steiner system with $k\geq t+1\geq4$ and a
$t$-transitive group, it leaves the affine binary plane designs,
projective subline designs, and the five Mathieu designs with parameters
\[
(4,5,11),\quad(5,6,12),\quad(3,6,22),\quad(4,7,23),\quad(5,8,24).
\]
The affine groups are $2^d:GL(d,2)$ and, when $d=4$, $2^4:A_7$.
In the subline case the blocks are the projective images of
$\mathbb F_q\cup\{\infty\}$ in $PG(1,q^e)$, $e\geq2$; the group
contains $PSL(2,q^e)$. We also use the upper bound by
$P\Gamma L(2,q^e)$ supplied by the proof of [K], not merely the lower
bound reproduced in [L, Theorem 10]. The classical actions in this
classification, including their identification with the native Mathieu
permutation groups in GAP, remain outside inputs. The actual designs,
incidences and subgroup calculations are checked separately below.
